\section{Breast Cancer Data Application}\label{sec:cancer-data-application}

The second empirical application considered a hospital-based cohort of
patients with breast cancer from Paraíba State, Brazil. Follow-up time was
represented by the recorded survival time, and death from breast cancer was
treated as the event of interest; observations without a recorded
breast-cancer death were treated as censored. In contrast with the melanoma application, the breast-cancer data
contained a substantially richer set of demographic, socioeconomic,
clinical, health-care access, lifestyle, and treatment-related information,
allowing the cure component to depend on a broader patient profile.

The predictor set included race or skin-colour classification, educational
level, income category, age at first consultation, age group, health
macroregion, health-system origin and referral pathway, and the hospital
source of the record. Clinical characteristics included primary tumour
site, histological type, and clinical stage, with stage grouped into
stages I--II, stages III--IV, and a separate category for unavailable
staging information. Lifestyle and history variables comprised tobacco
use, alcohol use, and family history of cancer. Information on the first
treatment received at the hospital was represented through separate
indicators for surgery, radiotherapy, chemotherapy, hormone therapy,
immunotherapy, no treatment, other treatments, and unavailable treatment
information. Health-care timing was represented by the intervals from
first consultation to diagnosis, diagnosis to first treatment, and first
consultation to first treatment. In addition to their continuous versions,
these intervals were represented by indicators distinguishing delays of up
to 90 days from delays exceeding 90 days.

Preprocessing was estimated exclusively from each training sample and then
applied unchanged to the corresponding validation or test observations.
Continuous age and delay variables were transformed using the Yeo--Johnson
power transformation\cite{yeo2000new} and standardized. Explicit missingness indicators were
retained for treatment-delay variables, while several categorical predictors
included dedicated levels for unavailable information. Because primary
tumour site and histological type contained more complex categorical
structures, they were represented using cross-validated survival-target
encodings based on restricted mean survival time, with smoothing applied
within the training data. The resulting representations were carried
unchanged to held-out observations, thereby preventing outcome information
from the validation or test samples from entering the preprocessing stage.

Several prespecified interaction terms were also included to allow clinically
relevant associations to vary across patient subgroups. These comprised
interactions between clinical stage and first treatment, between stage and
hospital source, and between stage and educational level. The continuous
consultation-to-diagnosis delay was additionally interacted with stage and
with hormone-therapy and immunotherapy indicators. These interaction terms,
together with the original transformed continuous variables, survival-target
encodings, and categorical indicators, formed the predictor matrix used by
all competing models.

Predictive performance was assessed using 50 event-stratified paired outer
holdouts, with approximately 70\% of the cohort assigned to training and
30\% to testing in each split. The same outer split indices were used for
all fitted methods. Hyperparameter selection was performed exclusively
within the outer training sample using four inner folds and the IPCW
integrated Brier score, together with the prediction-horizon rule defined
in Section~\ref{sec:prediction-evaluation}. The comparison comprised GG-KM
models under Poisson and negative-binomial latent-count laws, with linear,
radial-basis, Laplacian, exponential, Cauchy, polynomial, and
wavelet kernels, as well as spline Poisson and negative-binomial
models and the classical mixture-cure and promotion-time comparators.
All kernel and spline hyperparameters were selected from prespecified
grids.

\GGKMTable{sections/generated/breast-cancer-performance-table}

Table~\ref{tab:breast-cancer-performance} reports split-level means and
variation. The spline models under both Poisson and negative-binomial
latent-count laws achieved the lowest mean IPCW integrated Brier score,
0.0877 with split standard deviation 0.0019, followed by the mixture-cure
and promotion-time models, both with mean IBS 0.0883. The mixture-cure
model attained the highest mean cure AUC overall, 0.7851, followed by the
promotion-time model at 0.7775 and the spline models at 0.7707. Among the
GG-KM specifications, the negative-binomial linear kernel had the lowest
mean IBS, 0.0897, whereas the exponential kernel produced the highest mean
cure AUC, 0.7124, under both latent-count laws. Across the 900 planned
split-model fits, 873 completed successfully and 871 satisfied the explicit
convergence criteria.

Figure~\ref{fig:breast-cancer-performance} summarizes out-of-sample survival
calibration, prediction error, and cure discrimination. The split-averaged
predicted survival curves closely followed the Kaplan--Meier estimate over
the common prediction horizon, with relatively small differences among the
fitted methods. The spline Poisson and negative-binomial models showed the
lowest IPCW integrated Brier scores, whereas the mixture-cure model achieved
the highest cure-discrimination performance. The GG-KM models with the
exponential kernel were less competitive in overall prediction error than
the spline and classical cure models, but retained substantially better cure
discrimination than several alternative kernel specifications.

\GGKMWideFigure{breast_cancer_predictive_assessment_ibs_auc.pdf}
{Predictive assessment in 50 paired repeated holdouts of the breast cancer
cohort. Panel A compares the full-cohort Kaplan--Meier curve with
split-averaged out-of-sample survival predictions over the common supported
horizon. Panel B shows split-level IPCW integrated Brier scores and Panel C
shows the corresponding cure-discrimination area under the curve
($\mathrm{AUC}_{\mathrm{cure}}$). Each point represents one outer split, and
lower IBS values indicate better prediction whereas larger cure AUC values
indicate better discrimination.}
{fig:breast-cancer-performance}
{Three-panel predictive assessment for the breast cancer application. Panel A
compares the Kaplan--Meier survival curve with mean out-of-sample survival
predictions from the mixture-cure, promotion-time, spline, and selected GG-KM
models. Panels B and C show the corresponding repeated-holdout distributions
of IPCW integrated Brier scores and cure AUC values.}

Figure~\ref{fig:breast-cancer-performance-ggkm-kernel} shows that kernel
choice had a marked influence on GG-KM performance. For prediction error,
the exponential and Cauchy kernels yielded the lowest mean IBS values among
the nonlinear GG-KM specifications, while the wavelet and RBF kernels were
less favorable. Cure discrimination varied more substantially across
kernels. The exponential kernel produced the highest mean cure AUC among the
GG-KM models, 0.7124 under both count laws, whereas the wavelet kernel showed
considerably weaker discrimination.

% The linear Poisson model was particularly
% poor for cure discrimination, with mean AUC close to 0.5, while its
% negative-binomial counterpart performed substantially better. For most of
% the RBF, Laplacian, exponential, Cauchy, and wavelet specifications, however,
% the Poisson and negative-binomial results were very similar.

\GGKMWideFigure{breast_cancer_ggkm_kernel_comparison_ibs_auc.pdf}
{Kernel-specific predictive performance of the GG-KM models in 50 paired
repeated holdouts of the breast cancer cohort. Panel A shows the split-level
IPCW integrated Brier score and Panel B the cure-discrimination area under
the curve ($\mathrm{AUC}_{\mathrm{cure}}$) across the seven candidate
kernels. Within each kernel, Poisson and negative-binomial GG-KM results are
shown side by side, with color and hatching distinguishing the two
latent-count laws. Individual points represent outer splits and boxplots
summarize their distributions.}
{fig:breast-cancer-performance-ggkm-kernel}
{Two-panel comparison of GG-KM kernels in the breast cancer application.
Panel A displays split-level IPCW integrated Brier scores and Panel B cure
AUC values for the linear, RBF, Laplacian, exponential, Cauchy, polynomial,
and wavelet kernels. Poisson and negative-binomial specifications are shown
side by side for each kernel.}
